\begin{abstract}
As large language models (LLMs) become increasingly capable of generating human-like text, the ability to detect AI-generated content has emerged as a critical challenge for maintaining trust in digital communication. In parallel, understanding methods that reduce detectability is essential for developing robust detection systems and enabling legitimate use cases requiring natural-sounding AI text. We introduce a multi-track leaderboard framework for systematically benchmarking AI text undetectability methods, organizing techniques by their intervention point: inference-time prompt modifications, post-editing paraphrasing, fine-tuning, and pretraining. We implement and evaluate two tracks---inference and post-editing---using an ensemble detector combining perplexity-based and burstiness-based detection. Our experiments on GPT-4o-mini generated text reveal that inference-time modifications (\textit{e.g.}, human-style prompts) achieve 70\% evasion rates at 5\% false positive rate while preserving complete semantic content. Post-editing methods match this evasion performance but incur significant quality degradation (41\% word overlap). Surprisingly, personalization strategies that attempt to add human-like elements actually \textit{increase} detectability, suggesting that LLM-style personalization patterns are recognizable signatures. Our findings demonstrate that prompt engineering offers superior quality-evasion trade-offs compared to post-hoc paraphrasing, with implications for both AI writing assistant design and detection system development.
\end{abstract}
